% ================================================================== Chapitre 2
\chapter{Enthalpie libre et potentiel chimique}
\label{ch:potentiel}

En thermodynamique physique, les transformations étudiées ont lieu le plus souvent à température et pression fixées (réacteur ouvert à l'air, bain thermostaté). La fonction d'état adaptée à ces conditions est l'\emph{enthalpie libre} $G$. Pour un système dont la composition change, on lui associe une grandeur par constituant : le \emph{potentiel chimique}.

\section{La fonction enthalpie libre}

\begin{definition}[enthalpie libre]
L'\textbf{enthalpie libre} (ou énergie de Gibbs) d'un système est la fonction d'état
\[ G = H - TS = U + PV - TS \]
où $U$ est l'énergie interne, $H = U + PV$ l'enthalpie et $S$ l'entropie. C'est une grandeur \emph{extensive}, en joules.
\end{definition}

\begin{resultat}{Identité thermodynamique -- système fermé de composition fixée}
\[ \dd G = V\,\dd P - S\,\dd T \]
Les variables naturelles de $G$ sont $T$ et $P$ : à $T$ et $P$ fixées, $G$ d'un système de composition fixée ne varie pas.
\end{resultat}

\begin{demo}
On différentie $G = U + PV - TS$ : $\dd G = \dd U + P\,\dd V + V\,\dd P - T\,\dd S - S\,\dd T$. Pour une transformation réversible élémentaire, $\dd U = \delta W + \delta Q = -P\,\dd V + T\,\dd S$, d'où
\[ \dd G = (-P\,\dd V + T\,\dd S) + P\,\dd V + V\,\dd P - T\,\dd S - S\,\dd T = V\,\dd P - S\,\dd T. \]
Comme $G$ est une fonction d'état, cette relation entre états voisins reste vraie pour toute transformation.
\end{demo}

\begin{pieges}
\begin{itemize}
  \item $G = U \mathbf{+} PV - TS$ : le signe devant $PV$ est \emph{plus} (car $H = U + PV$). Une erreur de signe ici rend fausse toute la suite du calcul de $\dd G$.
  \item $\dd G = V\,\dd P - S\,\dd T$ suppose le système \emph{fermé} et \emph{de composition fixée}. Dès qu'une réaction a lieu, il faut ajouter les termes $\mu_i\,\dd n_i$ (section suivante).
\end{itemize}
\end{pieges}

\section{Système de composition variable : potentiel chimique}

Dès qu'une transformation chimique a lieu, les quantités de matière $n_i$ varient : $G$ devient une fonction $G(T, P, n_1, n_2, \dots)$.

\begin{definition}[potentiel chimique]
Le \textbf{potentiel chimique} du constituant $\mathrm{A}_i$ dans un mélange est
\[ \mu_i = \left(\frac{\partial G}{\partial n_i}\right)_{T,\,P,\,n_{j\neq i}} \qquad\text{en }\mathrm{J\cdot mol^{-1}}. \]
C'est une grandeur \emph{intensive} : la variation d'enthalpie libre du système quand on y ajoute une mole de $\mathrm{A}_i$ en gardant $T$, $P$ et les autres quantités constantes (à la limite d'un ajout infinitésimal).
\end{definition}

\begin{resultat}{Différentielle et expression de $G$ pour un mélange}
\[ \dd G = V\,\dd P - S\,\dd T + \sum_i \mu_i\,\dd n_i \qquad\qquad G = \sum_i n_i\,\mu_i \]
\end{resultat}

La seconde relation (\emph{identité d'Euler}) est admise : elle traduit l'extensivité de $G$. On peut s'en convaincre en « remplissant » le système à $T$, $P$ et composition fixées (donc à $\mu_i$ constants) : en ajoutant les constituants dans les bonnes proportions, chaque $\dd n_i$ apporte $\mu_i\,\dd n_i$ et l'on obtient $G = \sum_i \mu_i n_i$.

\begin{exemple}[corps pur seul]
Pour un corps pur, $G = n\,\mu$, donc $\mu = G/n = G_m$ : le potentiel chimique d'un corps pur est son \textbf{enthalpie libre molaire}. Pour une mole, l'identité thermodynamique devient $\dd\mu = V_m\,\dd P - S_m\,\dd T$ avec $G_m = H_m - TS_m$.
\end{exemple}

\section{Expression du potentiel chimique}

\begin{resultat}{Potentiel chimique en fonction de l'activité (admis)}
\[ \boxed{\mu_i(T, P, \text{composition}) = \mu_i^\circ(T) + RT\ln a_i} \]
$\mu_i^\circ(T)$ est le \textbf{potentiel chimique standard} de $\mathrm{A}_i$ : sa valeur quand $a_i = 1$, c'est-à-dire quand $\mathrm{A}_i$ est dans son état standard. Il ne dépend que de $T$.
\end{resultat}

Avec le tableau des activités du chapitre précédent :
\[ \mu_i^{\text{gaz}} = \mu_i^\circ + RT\ln\frac{P_i}{P^\circ}, \qquad \mu_i^{\text{soluté}} = \mu_i^\circ + RT\ln\frac{c_i}{c^\circ}, \qquad \mu^{\text{solide ou liquide pur}} = \mu^\circ. \]

\begin{horsprog}[origine du terme $RT\ln a_i$]
Pour un gaz parfait pur, $\dd\mu = V_m\,\dd P = \frac{RT}{P}\,\dd P$ à $T$ fixée ; en intégrant de $P^\circ$ à $P$, on obtient $\mu = \mu^\circ + RT\ln(P/P^\circ)$. Les autres expressions s'obtiennent par analogie (modèle de la solution idéale). Seul le résultat est au programme.
\end{horsprog}

\section{États standard}

\begin{definition}[état standard d'un constituant à la température $T$]
L'\textbf{état standard} d'un constituant est un état de référence, choisi par convention à la \emph{même température $T$} que le système et sous la \textbf{pression standard} $P^\circ = \SI{1}{bar}$ :
\begin{itemize}
  \item \textbf{gaz} : le constituant pur, se comportant comme un gaz parfait, sous $P^\circ$ ;
  \item \textbf{solide ou liquide} : le constituant pur, dans le même état physique, sous $P^\circ$ ;
  \item \textbf{soluté} : le soluté à la concentration $c^\circ = \SI{1}{mol.L^{-1}}$, se comportant comme en solution infiniment diluée ;
  \item \textbf{solvant} : le solvant pur, sous $P^\circ$.
\end{itemize}
\end{definition}

\begin{pieges}
\begin{itemize}
  \item Il n'y a \emph{pas} de température standard : il existe un état standard \emph{à chaque température}. Les tables sont souvent données à $\SI{298}{K}$ par commodité, mais « standard » ne veut pas dire « $\SI{25}{\celsius}$ ».
  \item L'état standard peut être \emph{fictif} : un gaz réel pur n'est jamais exactement parfait ; un soluté à $\SI{1}{mol.L^{-1}}$ n'est jamais infiniment dilué. C'est une référence de calcul, pas un état réel.
  \item Pour \ce{H+} en solution, l'état standard correspond à $[\ce{H+}] = c^\circ$, donc à $\mathrm{pH} = 0$ : c'est pourquoi les potentiels standard des couples faisant intervenir \ce{H+} sont ceux mesurés à $\mathrm{pH} = 0$.
\end{itemize}
\end{pieges}

\begin{definition}[état standard de référence d'un élément]
L'\textbf{état standard de référence} (ESR) d'un élément chimique à la température $T$ est l'état standard du \emph{corps simple} correspondant, dans l'état physique \emph{le plus stable} à cette température (sous $P^\circ$).

\noindent\emph{Exceptions conventionnelles} : pour H, N, O, F et Cl, l'ESR est le gaz diatomique (\ce{H2}, \ce{N2}, \ce{O2}, \ce{F2}, \ce{Cl2}) \emph{à toute température} ; pour le carbone, c'est le graphite.
\end{definition}

\begin{exemple}[ESR à différentes températures]
\begin{itemize}
  \item Oxygène : \ce{O2(g)} et non \ce{O3(g)} (l'ozone est un corps simple moins stable).
  \item Brome : \ce{Br2(l)} à $\SI{298}{K}$, mais \ce{Br2(g)} au-dessus de sa température d'ébullition $\SI{332}{K}$.
  \item Fer : \ce{Fe(s)} jusqu'à $\SI{1811}{K}$, puis \ce{Fe(l)}.
  \item Carbone : graphite et non diamant.
\end{itemize}
\end{exemple}

\begin{exemple}[graphite ou diamant ?]
À $\SI{298}{K}$, $\mu^\circ(\text{C}_{\text{diamant}}) - \mu^\circ(\text{C}_{\text{graphite}}) = +\SI{2,9}{kJ.mol^{-1}} > 0$. Les deux étant des solides purs ($a = 1$), $\mu = \mu^\circ$ : le diamant a un potentiel chimique plus élevé. La transformation diamant $\to$ graphite est donc thermodynamiquement favorisée (on le démontre au chapitre~\ref{ch:evolution}), mais elle est infiniment lente à température ambiante : le diamant est \emph{métastable}. La thermodynamique dit ce qui \emph{peut} se produire, la cinétique dit \emph{à quelle vitesse}.
\end{exemple}
